\documentclass[10pt]{article}

\usepackage[utf8]{inputenc}

\usepackage[T1]{fontenc}

\usepackage{amsmath}

\usepackage{amsfonts}

\usepackage{amssymb}

\usepackage[version=4]{mhchem}

\usepackage{stmaryrd}

\usepackage{bbold}

\usepackage{mathrsfs}

\usepackage{graphicx}

\usepackage[export]{adjustbox}

\graphicspath{ {./images/} }



\begin{document}

оценок рассматриваемого типа. Однако для ряда задач невычислительного характера такого типа оценки известны (см. Алгоритма сложность).



В случае, когда исследование проблемы О.в. а. имеет своей целью решение задач на ЭВМ, проблема О. в. а. приобретает дополнительные оттенки, связанные с устойчивостью алгоритма к вычислительной погрешности, с ограничением на объем разных видов используемой памяти (см. Вычислительный алгоритм). Выше задача О. в. а. рассматривалась как задача О. в. а. на классах задач. Супественный интерес для практики представляет задача О. в. а. на конкретной задаче (см. [10], [16]). Постановка проблемы оптимизации (см. [16]) заключается в следующем. Дифференциальное уравнение интегрируется методом Рунге - Кутта с переменным шагом. Производится оценка главного члена оценки погрешности. Затем әта оценка оптимизируется по распределению узлов интегрирования (при общем заданном числе узлов). Такой подход к проблеме оптимизации оказал существенное влияние на развитие теории и практики активных алгоритмов численного интегрирования.



Лит.: [1] Н и к о л ь с к и й С. М., Квадратурные формулы, 3 изд., М., 1979; [2] К о л м о г о р о в А. Н., "Докл. АН СССР", 1956 , т. 108, N 3 , $385-88$; [3] К о л м о г о р о в А. Н., Т их о м и $\mathrm{p}$ ов B. М., "Успехи матем. наук», 1959, т. 14, в. 2, т. Заулирования, М., 1959 ; [5] Б а 6 у у $ш$ а И., С о б о л е в С. Л., табулирования, М., 1959 ; [5] Б а б у Ш К а И., С о б о л е в С. Л., "Aplikace mat.", 1965, t. 10, s. 96-129; [6] Б а х в а Л о в H. С., ный конгеесс математиков в Ницце. 1970. Доклады советских математиков, М., 1972 , с. $27-33$; [8] е г о ж е, в кн.: Численные методы решения дифференциальных и интегральных уравнений и квадратурные формулы, М., 1964, с. 3-63; [9] е г о ж е, «Матем. заметки», 1972 , т. 12 , в. 6, с. 655-64; [10] е г о ж е, Численные методы, 2 изд., М., 1975 ; [11] е г о ж е, «Ж. вычисЧисленные методы, 2 изд., М., 1975; [11] е го Ж е, «Ж. вычис-

лит. матем. и матем. физ.», 1971, т. 11, № 4, с. $1014-18$; [12] Лит. матем. и матем. Физ.», 1971, т. 11, № 4, с. 1014-18; [12] М., 1967; [13] В а си л ь е в Ф. П., Численные методы решения экстремальных задач, М., 1980; [14] О г а н е с я н Л. А., Р ух о в е ц Л. А., Вариационно-разностные методы решения элпиптических уравнений, Ер., 1979; [15] Т и х о м и р о в В. М., Некоторые вопросы теории приближений, М., 1976; [16] Т и х он о в А. Н., Г о р б у н о в А. Д., “Ж. вычислит. матем. и матем. физ.», 1964, т. 4, № 2, с. 232 - 41 . H. С. Бахвалов



ОПЦИОНАЛЬНАЯ б-АЛГЕБРА - наименышая $\sigma$-алгебра $\mathcal{G}=G(\mathbb{F})$ множеств из $\Omega \times R_{+}=\{(\omega, t): \omega \in \Omega$, $t \geqslant 0\}$, порожденная всеми отображениями $(\omega, t) \sim$ $f(\omega, t)$ множества $\Omega \times R_{+}$в $R$, к-рые (для каждого фиксированного $\omega \in \Omega$ ) являются (по $t$ ) непрерывными справа, имеют пределы слева и $\mathbb{F}$-согласованы с неубывающим семейством $\mathbb{F}=\left(F_{t}\right)_{t \geqslant 0}$ под-б-алгебр $F_{t} \subseteq F$, $t \geqslant 0$, где $(\Omega, F)$ - измеримое пространство. 0 . $\sigma-$ совпадает с наименьшей $\sigma$-алгеброй, порожденной стохастич. интервалами $\llbracket 0, \tau \rrbracket=\{(\omega, t): 0 \leqslant t<\tau(\omega)\}$, где $\tau=\tau(\omega)$ - марковские моменты (относительно $\mathbb{F}=$ $\left.=\left(F_{t}\right)_{t \geqslant 0}\right)$. Между опциональными и предсказуемыми б-алгебрами имеет место соотношение $\mathscr{P}(\mathbb{F}) \sqsubseteq \mathcal{G}(\mathbb{F})$.



Лит.: [1] Д е л л а ш е р и К., Емкости и случайные процессы, пер. с франц., М., стохастический процесс $X=\left(X_{t}(\omega), F_{t}\right)_{t \geqslant 0}$, являющийся измеримым (как отображение $(\omega, t) \sim \rightarrow(\omega, t)=$ $\left.=X_{t}(\omega)\right)$ относительно опциональной $\sigma$-алгебры $G=$ $=G(\mathbb{F})$.



A. Н. Ширлев.



ОРБИТ МЕТОД - метод изучения унитарных представлений групп Ли. С помощью О. м. была построена теория унитарных представлений нилыпотентных групп Ли, а также указана возможность его применения к другим группам [1].



О. м. основан на следующем «экспериментальном» факте: существует глубокая связь между унитарными неприводимыми представлениями группы Ли $G$ и орбитами әтой группы в коприсоединенном представлении. Решение основных задач теории представлений с помощью О. м. осуществляется следующим образом (см. [2]). I. К овструкция п классификапия



\includegraphics[max width=\textwidth, center]{2023_07_08_9b382d9aca2441937daag-1}

ле н и й. Пусть $\Omega$ - орбита действительной группы Ли $G$ в коприсоединенном представлении, $F$ - точка этой орбиты (являющаяся линейным функционалом на алгебре Ли q группы $G), G(F)$ - стабилизатор точки $F, \mathfrak{g}(F)$ - алгебра Ли группы $G(F)$. П о л я р и з ац и е й точки $F$ наз. комплексная подалгебра $\mathfrak{h}$ в $\mathfrak{g}_{\mathbb{C}}$, где $\mathfrak{g}_{\mathbb{C}}-$ колплексибикация алгебры $Л и \mathfrak{g}$, обладающая свойствами:



\begin{enumerate}

  \item $\operatorname{dim}_{\mathbb{C}} \mathfrak{h}=\operatorname{dim} \mathfrak{g}-\frac{1}{2} \operatorname{dim} \Omega$



  \item $[\mathfrak{h}, \mathfrak{h}]$ содержится в ядре функционала $F$ на $\mathfrak{g}$; 3) $\mathfrak{h}$ инвариантна относительно Ad $G(F)$.



\end{enumerate}



Пусть $H^{0}=\exp (\mathfrak{h} \cap \mathfrak{q})$ и $H=G(F) \cdot H^{0}$. Поля ризация $\mathfrak{h}$ наз. де й с т в и те ль н о й, если $\mathfrak{h}=\overline{\mathfrak{h}}$, и ч и с т о к о мп л е к с н о й, если $\mathfrak{h} \cap \overline{\mathfrak{h}}=\mathfrak{g}(F)$. Функционал $F$ определяет характер (одномерное унитарное представление) $\chi_{F}^{0}$ группы $H^{0}$ по формуле



$\exp X \longrightarrow \exp 2 \pi i\langle F, X\rangle$.



Пусть $\chi_{F}^{0}$ продолжается до характера $\chi_{F}$ группы $H$. Если $\mathfrak{h}$ - действительная поляризация, то пусть $T_{F, \mathfrak{h}, \chi_{F}}$ - индуцированное характером $\chi_{F}$ подгруппы $H$ представление группы $G$ (см. Индуцированное представление). Если $\mathfrak{h}$ - чисто комплексная поляризация, то пусть $T_{F, \mathfrak{h}, \chi_{F}}$ - голоморфно индуцированное представление, действующее в пространстве голоморфных функций на $G / H$.



П е р в а о с но в ная г и по те том, что представление $T_{F, \mathfrak{l}}, \chi_{F}$ неприводимо и его класс эквивалентности зависит только от орбиты $\Omega$ и от выбора продолжения $\chi_{F}$ характера $\chi_{F}^{0}$. Эта гипотеза доказана для нилыпотентных групп [1] и для разрешимых групп Ли [5]. Для нек-рых орбит простой особой группы $G_{2}$ гипотеза неверна [7]. Возможность продолжения и степень его неодновначности зависят от топологич. свойств орбиты: препятствием к продолжению служат двумерные когомологии орбиты, а в качестве параметра, нумерующего различные продолжения, можно взять одномерные когомологии орбиты. Более точно, пусть $B_{\Omega}$ - каноническая 2-форма на орбите $\Omega$. Для существования продолжения необходимо и достаточно, чтобы форма $B_{\Omega}$ принадлежала целочисленному классу (т. е. интеграл ее по любому двумерному циклу был целым числом); если это условие выполнено, то множество продолжений параметризуется характерами фундаментальной группы орбиты.



В торая ос новная ги потеза состоит в том, что указанным способом получаются все унитарные неприводимые представления рассматриваемой группы G. Единственным (1983) противоречащим әтой гипотезе примером являются т. н. дополнительные серии представлений полупростых групп Ли.



II. Ф у нк то р и а льны е с в о й с т в а с о о тветствия ме жу ор б итам и и предс т а в л е н и я м и. Значительное место в теории представлений занимают вопросы о разложении на неприводимые компоненты представления, получаемого ограничением на подгруппу $H$ неприводимого представления группы $G$ и индуцированием с помощью неприводимого представления подгруппы $H \subset G$. O. м. дает ответ на әти вопросы в терминах естественной проекции $p: \mathfrak{g}^{*} \rightarrow \mathfrak{h}^{*}(*$ означает переход $\boldsymbol{\kappa}$ сопряженному пространству; проекция $p$ состоит в ограничении функционала с $\mathfrak{g}$ на $\mathfrak{h})$ ). А именно, густь $G$ - экспоненциальная группа Ји (для таких групп соответствие между орбитами и представлениями взаимно однозначно).





\end{document}